\chapter*{Özet}
\addcontentsline{toc}{chapter}{Özet}

Bu projede; taksi gibi araç içi ortamlarda sürücü güvenliğini artırmaya
yönelik, gömülü sistem tabanlı bir yüz tanıma çözümü tasarlanmış ve
prototip düzeyinde gerçekleştirilmiştir. Sistem, arka koltuğa yerleştirilen
bir ESP32-CAM modülü üzerinden 2 saniye boyunca 5 kare/saniye hızında
toplam 10 karelik bir görüntü öbeği yakalar ve Wi-Fi üzerinden bulutta
çalıştırılan tanıma servisine iletir. Bulut servisi açık kaynak
InsightFace kütüphanesini (buffalo\_l, ArcFace-R100) kullanarak çoklu
kare için kalite filtresi uygular, geçerli gömme vektörlerinin
centroidini hesaplar ve aranan kişi veri tabanı ile kosinüs benzerliği
ölçer. Sonuç, UART üzerinden STM32 NUCLEO-F767ZI tabanlı olay
yöneticisine iletilir; STM32 LED ve buzzer ile sürücüye geri bildirim
verir ve HM-10 BLE modülü üzerinden eşli Android telefona iletir.
Telefondaki uygulama, BLE bildirimini alarak otomatik biçimde
\textbf{155 Polis İmdat} hattını arar. Tasarım modülerdir; donanım yan~ürünü
olarak elde edilen yerel ön~işleme, bulut tanıma ve BLE tetikleme
katmanları birbirinden bağımsız geliştirilebilir niteliktedir. Çalışma,
düşük maliyetli (donanım eki $<$~500~TL), KVKK uyumlu ve filo araçlarına
genişletilebilir bir referans mimari sunmaktadır.

\noindent\textbf{Anahtar Kelimeler:} yüz tanıma, ArcFace, gömülü sistem,
ESP32-CAM, STM32, BLE, IoT güvenlik, taksi güvenliği, çoklu-kare
agregasyon, pasif canlılık.

\vspace{1cm}

\chapter*{Abstract}
\addcontentsline{toc}{chapter}{Abstract}

This thesis presents the design and prototype implementation of an
embedded face-recognition system aimed at improving driver safety in
in-vehicle environments such as taxis. An ESP32-CAM module mounted toward
the rear seat captures a 10-frame burst at 5\,FPS over 2 seconds upon the
driver's scan command and uploads it over Wi-Fi to a cloud recognition
service. The service, built on top of the open-source InsightFace
library (buffalo\_l, ArcFace-R100) and hosted on AWS EC2, applies a
per-frame quality gate, computes the centroid of valid embeddings, and
performs cosine similarity matching against a wanted-suspect database.
The result is delivered over UART to an STM32 NUCLEO-F767ZI based event
manager that drives the driver-facing LEDs and buzzer and forwards the
event to a paired Android phone over an HM-10 Bluetooth Low Energy link.
The Android application listens for BLE notifications and automatically
dials the national emergency number \textbf{155}. The architecture is
modular: on-device pre-processing, cloud-side recognition and BLE-driven
actuation can be developed independently. The result is a low-cost
(under \euro\,15 in extra hardware), GDPR-aware, fleet-extensible
reference architecture.

\noindent\textbf{Keywords:} face recognition, ArcFace, embedded systems,
ESP32-CAM, STM32, BLE, IoT security, taxi safety, multi-frame aggregation,
passive liveness.

\cleardoublepage
